\chapter{VSIDS}

Now that we've finished implementing CDCL, these next chapters will be focused 
on other algorithms and data structures used to improve the performance of CDCL.
We will start with one of the easier methods of improving CDCL: VSIDS.

Variable State Independent Decaying Sum (VSIDS) is a type of variable assignment 
heuristic, basically a loosely defined rule, that is used in modern CDCL solvers.
It is \textit{the} best variable assignment/branching heuristic with many variations.
In this chapter, we follow this paper
\begin{center}
    \url{https://arxiv.org/pdf/1506.08905}
\end{center}
which gives a much more in depth and theoretical explanation of how and why they work.
Four our purposes, we'll understand at a higher level what VSIDS are as well as how 
and why they work. 

Currently, our variable assignment choosing heuristic has been ``find first Unassigned
variable, and assign it True.'' With VSIDS, this becomes ``find the most active Unassigned
variable, and assign it True.'' The first step in VSIDS is to add in an array of
\textbf{activity scores} which rank ``how active'' a variable is. There are two schemes 
for ranking activity score, which is ranking each literal (positive and negative variables),
or just the variables (absolute value of the literal). For our purposes, we will rank just 
the variables. 

The next step is after clause resolution, add to each variable's score if it was involved 
in clause resolution. Again, there are multiple ways to do this. Some implementations choose 
only the variables in the learned clause, some implementations choose every variable involved 
in the resolution algorithm, even if they are resolved away, and there are more implementations 
than just those two. For our purposes, we will choose bumping the scores of only the variables 
in the learned clause.

The last step, which doesn't happen every iteration, is to multiply each score by a decay factor 
\footnote{Not to be confused with Deadpool's ``dying factor."}
after a certain number of conflicts. The idea behind this is to even the playing field for 
each variable. If a variable accumulates a giant activity score, obviously it will be the first 
pick if it's valid. However, activity scores don't represent how recent the activity score bump 
occurred, so if that variable stopped being active, it would still be the first pick despite 
not actually being active relative to the learned clauses. This decay lets the variables that 
are starting to become more active have a better chance at competing with the potentially 
inactive, but score-dominating variables. As it is a decaying factor, this value is chosen 
to be greater than 0, but less than 1, so that the scores shrink to some percentage.

\section{cVSIDS}

We're going to start out with implementing the most basic form of VSIDS and work our way 
up to some better versions. We'll start with the ``textbook'' form, which was outlined in 
the section's introduction. This form of VSIDS is called \textbf{cVSIDS} by the paper we're 
using. We first add in our \textit{activity} array into our solver 
object with a counter for the number of conflicts.
\begin{lstlisting}
    public interface:
        SATSolver(string input)
        bool solveCNF()

    private interface:
        ...
        double[] activity
        int numConflicts
        ...
\end{lstlisting}
Then our solve code becomes as follows:

\begin{algorithm}[H]
\caption{cVSIDS Solving Method}\label{alg:36}
\begin{algorithmic}[1]
\Procedure{solve}{}

\State \Call{setup}{}

\While{true}
    \If{trailHead = size of trail}
        \State maxScore $\gets$ -1.0
        \State firstUnassigned $\gets$ -1
        \For{i $\gets$ 1 up to numVars}
            \If{assignment[i] = Unassigned}
                \If{maxScore < activity[i]}
                    \State firstUnassigned $\gets$ i
                    \State maxScore $\gets$ activity[i]
                \EndIf
            \EndIf
        \EndFor

        \State $\Call{assign}{firstUnassigned}$
    \EndIf

    \State propResult $\gets \Call{propagate}{assignment}$
    \If{propResult failed}
        \State numConflicts $\gets$ numConflicts + 1

        \State \Call{resolveConflict}{conflict}

        \State $\Call{backjump}{backJumpTo}$
        \State \Comment{apply decay factor}
        \If{numConflicts = MAX\_CONFLICTS}
            \For{i $\gets$ 1 up to numVars}
                \State activity[i] $\gets$ activity[i] * DECAY
            \EndFor
            \State numConflicts $\gets$ 0
        \EndIf

        \State continue
    \EndIf
\EndWhile

\EndProcedure
\end{algorithmic}
\end{algorithm}

\begin{algorithm}[H]
\caption{cVSIDS First UIP resolution}\label{alg:37}
\begin{algorithmic}[1]
\Procedure{resolveConflict}{conflict[1..n]}

\State ...
\While{seenIdx is not empty}
    \State lIdx $\gets$ seenIdx.last 
    \State activity[lIdx] $\gets$ activity[lIdx] + 1
    \State ...
\EndWhile

\EndProcedure
\end{algorithmic}
\end{algorithm}

The pseudocode is very straightforward, so I won't bother to explain it much. Just 
so you know, if you choose to initialize \textit{maxScore} with a value of 0.0 in the 
assignment code, use \verb|>=| instead of \verb|>| because the activity scores start 
out at 0.0, so not using \verb|>=| will mess with unsatisfiable problems. In our 
code to apply the decay factor, I use ``X'' and ``Y'' to represent the number of 
conflicts and the decay factor respectively. The reason why these aren't set values,
but rather placeholders is because there aren't two numbers that perform the best for
\textit{all} equations, so they should be adjusted programmatically or manually set 
experimentally. However, I just decided to choose two random numbers, so I used 
\textbf{X = 10} and \textbf{Y = 0.8}. But results for different X and Y will be 
shown as well.

The results were as follows:
\noindent
uf20.91: \textbf{279ms}

\noindent
UUF50.218.1000: \textbf{1841ms}

We didn't get a big speedup for satisfiable problems, but we 
did shave off about half a second for the unsatisfiable problems, 
which I think is pretty darn good.

\begin{FVerbatim}
On all 1000 equations in uf20-91:

cVSIDS (X = 1, Y = 0.5):    260ms (~0.26 seconds).
cVSIDS (X = 1, Y = 0.8):    266ms (~0.27 seconds).
cVSIDS (X = 10, Y = 0.5):   270ms (~0.27 seconds).
cVSIDS (X = 10, Y = 0.8):   279ms (~0.29 seconds).
Improved CDCL:              284ms (~0.28 seconds).

On all 1000 equations in UUF50.218.1000:

cVSIDS (X = 1, Y = 0.5):    1,680ms (~1.68 seconds).
cVSIDS (X = 1, Y = 0.8):    1,705ms (~1.71 seconds).
cVSIDS (X = 10, Y = 0.5):   1,782ms (~1.78 seconds).
cVSIDS (X = 10, Y = 0.8):   1,841ms (~1.84 seconds).
Improved CDCL:              2,385ms (~2.39 seconds).
\end{FVerbatim}
\noindent
In the results, since we will be interating over many versions of VSIDS, 
I will leave out the older results and only compare against the previous 
versions of our solver. In the results, it appears to be that lower 
X and Y values lead to faster solving times, however, I encourage you to 
choose your own X and Y values and see what happens for yourself.

\section{mVSIDS}

\textbf{mVSIDS} is the second form of VSIDS introduced by the paper from 
the section introduction. In this type of VSIDS, we bump the values of 
every variable we resolve on during clause resolution as well as the values 
of the variables in the learned clause. The multiplicative decay step is 
applied the same as during cVSIDS. This means that we only need to change the
clause resolution code.

The algorithmic change is basically trivial:

\begin{algorithm}[H]
\caption{mVSIDS First UIP resolution}\label{alg:38}
\begin{algorithmic}[1]
\Procedure{resolveConflict}{conflict[1..n]}

\State ...

\For{each $entry$ in $trail$ from last} \Comment traverse the trail in reverse order
    \State idx $\gets$ entry.lit
    
    \If{seen[idx] = true}
        \State reason $\gets$ Clauses[entry.reasonIdx]
        \For{each $var$ in $reason$}
            \State ...
        \EndFor

        \State ...
        \State activity[idx] $\gets$ activity[idx] + 1
    \EndIf

    \If{levelCount = 1}
        \State break
    \EndIf
\EndFor

\While{seenIdx is not empty}
    \State lIdx $\gets$ seenIdx.last 
    \State activity[lIdx] $\gets$ activity[lIdx] + 1
    \State ...
\EndWhile

\EndProcedure
\end{algorithmic}
\end{algorithm}

\noindent 
Here are the results using \textbf{X = 10} and \textbf{Y = 0.8}:

\begin{FVerbatim}
On all 1000 equations in uf20-91:

mVSIDS (X = 1, Y = 0.5):    279ms (~0.28 seconds).
mVSIDS (X = 1, Y = 0.8):    277ms (~0.28 seconds).
mVSIDS (X = 10, Y = 0.5):   269ms (~0.27 seconds).
mVSIDS (X = 10, Y = 0.8):   284ms (~0.28 seconds).

cVSIDS (X = 1, Y = 0.5):    260ms (~0.26 seconds).
cVSIDS (X = 1, Y = 0.8):    266ms (~0.27 seconds).
cVSIDS (X = 10, Y = 0.5):   270ms (~0.27 seconds).
cVSIDS (X = 10, Y = 0.8):   279ms (~0.29 seconds).

Improved CDCL:              284ms (~0.28 seconds).

On all 1000 equations in UUF50.218.1000:

mVSIDS (X = 1, Y = 0.5):    1,532ms (~1.53 seconds).
mVSIDS (X = 1, Y = 0.8):    1,526ms (~1.53 seconds).
mVSIDS (X = 10, Y = 0.5):   1,420ms (~1.42 seconds).
mVSIDS (X = 10, Y = 0.8):   1,467ms (~1.47 seconds).

cVSIDS (X = 1, Y = 0.5):    1,680ms (~1.68 seconds).
cVSIDS (X = 1, Y = 0.8):    1,705ms (~1.71 seconds).
cVSIDS (X = 10, Y = 0.5):   1,782ms (~1.78 seconds).
cVSIDS (X = 10, Y = 0.8):   1,841ms (~1.84 seconds).

Improved CDCL:              2,385ms (~2.39 seconds).
\end{FVerbatim}
\noindent
So even with the worst mVSIDS time, it beats the best cVSIDS time
(per the benchmark, results may vary, these tests are very small and are not 
very representative of real-world SAT solver usage).

\section{EMA/Normalized VSIDS}

EMA stands for ``Exponential Moving Average'' and it's what allows us to 
mathematically decide which variables' activity scores should be decayed 
and by how much. If we go back to the earlier example of some variable gaining 
a really high activity score and then becoming inactive, but still dominating 
over the more recently active variables, this is what EMA helps with. EMA will 
decay high scoring, but actually inactive variables' scores while mostly 
preserving the actually active variables' scores. 

The way it does it is with this nifty formula:
$$s_n = (1 - f)\delta_n + f * s_{n - 1}$$

$s_n$ is the \textbf{current score}, $f$ is our \textbf{decay factor},
$\delta_n \in \{0, 1\}$ where a value of $0$ corresponds to a variable 
not being bumped, and a value of $1$ corresponds to a variable being bumped.
Let's walk through an example:

\begin{center}
    Suppose $f = 0.95$, and $s_0 = 0.0$. In the first conflict ($n = 1$),
    a variable was involved. Its score becomes:
    $$s_1 = (1 - 0.95)(1) + 0.95(0.0) = 0.05$$
    It's not much, but it's much more than 0, so this variable is considered 
    high ranking. Now suppose in the next conflict ($n = 2$), the variable was 
    not involved. Its score now becomes:
    $$s_2 = (1 - 0.95)(0) + 0.95(0.05) = 0.0475$$
    This is less than $0.05$, so now it's low ranking, even if the difference is small.
    Now suppose it's involved in the third conflict ($n = 3$):
    $$s_3 = (1 - 0.95)(1) + 0.95(0.0475) = 0.095125$$
    This is much larger than what it was before.
\end{center}

Although you probably have guessed already, we've been implementing VSIDS 
in a straightforward naive way by following the described process down to the letter. 
For sake of simplicity, for Normalized VSIDS, we'll also implement it naively.

Implementing normalized VSIDS isn't too different to how we're already implementing it,
except now, we update every score upon each conflict iteration. We'll need an additional 
data structure to keep track of the $\delta$ values. Since we're appending normalization 
on top of mVSIDS, we still update a $\delta$ value only when we resolve a variable or it 
appears in the learned clause. This means that we now delay the updating of our activity 
scores until after we construct the learned clause:

\begin{algorithm}[H]
\caption{cVSIDS Solving Method}\label{alg:39}
\begin{algorithmic}[1]
\Procedure{solve}{}

\State \Call{setup}{}

\While{true}
    \If{trailHead = size of trail}
        \State maxScore $\gets$ -1.0
        \State firstUnassigned $\gets$ -1
        \For{i $\gets$ 1 up to numVars}
            \If{assignment[i] = Unassigned}
                \If{maxScore < activity[i]}
                    \State firstUnassigned $\gets$ i
                    \State maxScore $\gets$ activity[i]
                \EndIf
            \EndIf
        \EndFor

        \State $\Call{assign}{firstUnassigned}$
    \EndIf

    \State propResult $\gets \Call{propagate}{assignment}$
    \If{propResult failed}
        \State numConflicts $\gets$ numConflicts + 1

        \State \Call{resolveConflict}{conflict}

        \State $\Call{backjump}{backJumpTo}$
        \State \Comment{apply decay factor}

        \For{i $\gets$ 1 up to numVars}
            \If{delta[i] = true} 
                \State activity[i] $\gets$ (1 - Y) + DECAY * activity[i]
            \Else
                \State activity[i] $\gets$ DECAY * activity[i]
            \EndIf
            \State delta[i] $\gets$ false
        \EndFor
        \State numConflicts $\gets$ 0

        \State continue
    \EndIf
\EndWhile

\EndProcedure
\end{algorithmic}
\end{algorithm}

\begin{algorithm}[H]
\caption{Normalized VSIDS First UIP resolution}\label{alg:40}
\begin{algorithmic}[1]
\Procedure{resolveConflict}{conflict[1..n]}

\State ...

\For{each $entry$ in $trail$ from last} \Comment traverse the trail in reverse order
    \State idx $\gets$ entry.lit
    
    \If{seen[idx] = true}
        \State reason $\gets$ Clauses[entry.reasonIdx]
        \For{each $var$ in $reason$}
            \State ...
        \EndFor

        \State ...
        \State delta[idx] $\gets$ delta[idx]
        \State activity[idx] $\gets$ activity[idx] + 1
    \EndIf

    \If{levelCount = 1}
        \State break
    \EndIf
\EndFor

\While{seenIdx is not empty}
    \State lIdx $\gets$ seenIdx.last 
    \State activity[lIdx] $\gets$ activity[lIdx] + 1
    \State delta[lIdx] $\gets$ true
    \State ...
\EndWhile

\EndProcedure
\end{algorithmic}
\end{algorithm}

\noindent
Notice how that there is no X value anymore because normalization happens upon 
every iteration instead of after a set number. Thankfully, our fastest times were 
when we updated activity scores after every iteration anyways, so we probably weren't 
missing some ``fastest'' time under a different X value.

Here are the results:
\begin{FVerbatim}
On all 1000 equations in uf20-91:

EMA (X = 1, Y = 0.95):      266ms (~0.27 seconds).
EMA (X = 1, Y = 0.8):       287ms (~0.29 seconds).
EMA (X = 1, Y = 0.5):       280ms (~0.28 seconds).

mVSIDS (X = 1, Y = 0.5):    279ms (~0.28 seconds).
mVSIDS (X = 1, Y = 0.8):    277ms (~0.28 seconds).
mVSIDS (X = 10, Y = 0.5):   269ms (~0.27 seconds).
mVSIDS (X = 10, Y = 0.8):   284ms (~0.28 seconds).

cVSIDS (X = 1, Y = 0.5):    260ms (~0.26 seconds).
cVSIDS (X = 1, Y = 0.8):    266ms (~0.27 seconds).
cVSIDS (X = 10, Y = 0.5):   270ms (~0.27 seconds).
cVSIDS (X = 10, Y = 0.8):   279ms (~0.29 seconds).

Improved CDCL:              284ms (~0.28 seconds).

On all 1000 equations in UUF50.218.1000:

EMA (X = 1, Y = 0.95):      1,504ms (~1.50 seconds).
EMA (X = 1, Y = 0.8):       1,497ms (~1.50 seconds).
EMA (X = 1, Y = 0.5):       1,621ms (~1.62 seconds).

mVSIDS (X = 1, Y = 0.5):    1,532ms (~1.53 seconds).
mVSIDS (X = 1, Y = 0.8):    1,526ms (~1.53 seconds).
mVSIDS (X = 10, Y = 0.5):   1,420ms (~1.42 seconds).
mVSIDS (X = 10, Y = 0.8):   1,467ms (~1.47 seconds).

cVSIDS (X = 1, Y = 0.5):    1,680ms (~1.68 seconds).
cVSIDS (X = 1, Y = 0.8):    1,705ms (~1.71 seconds).
cVSIDS (X = 10, Y = 0.5):   1,782ms (~1.78 seconds).
cVSIDS (X = 10, Y = 0.8):   1,841ms (~1.84 seconds).

Improved CDCL:              2,385ms (~2.39 seconds).
\end{FVerbatim}
\noindent
Here we see that EMA/normalizing the VSIDS doesn't make a huge contribution,
but it still makes a decent one for some values of Y.

\section{AMA/adaptVSIDS}

Adaptive Moving Average (AMA) VSIDS is an extension to EMA VSIDS where rather than 
having a fixed decay factor, we choose between two different decay rates depending on 
the state of our learned clause. The big idea behind making the decay factor variable 
instead of fixed is that our solver's ``productivity" isn't fixed. The learned clause 
affects this productivity in how high the quality of the learned clause is. 

To figure out the quality of a learned clause, we have to look at its \textbf{LBD}
which stands for \textbf{literal blocks distance}. LBD is defined as the number of 
distinct decision levels in the learned clause. A learned clause with a low LBD is considered to be higher 
quality than that of one with a high LBD. The paper doesn't explicitly say why this is the case,
\footnote{I might be illiterate and may have just missed the reason.}
but we can make a few guesses as to why. Considering that a low LBD means that there aren't 
many decision levels involved in the contradiction, so that means that the contradiction 
is more or less ``confined" to a smaller area of the search space. This indicates that 
the conflict has a more localized relationship with the current set of decisions. 
Such clauses have a higher chance to become unit or nearly unit, under a smaller 
amount of changes to the current assignment, better contraining future search more effectively.
Contrast this to high LBD, which has multiple decision levels involved in the contradiction, 
meaning that the contradiction spans a broader portion of the current search space, making 
it less likely it'll have an immediate effect on constraining future search.

This is supported by the fact that when we learn a clause with a low LBD, we choose a 
higher decay factor, so that variable activity scores shrink less, meaning we remember 
them more in the present. Likewise, when we learn a clause with a high LBD, we choose 
a lower decay factor, so variable activity scores shrink faster, making older activity 
lose its influence faster relative to the newer active variables.

The paper introduces this variable called \textit{lbdema}, which is defined as the 
exponential moving average of the LBD. Like our EMA VSIDS, this is calculated 
upon every iteration, and it's also a recurrence relation where the previous value affects 
what the next value will be. The formula is similar to the previous EMA formula:
$${lbdema}_n = (1 - \alpha) * {LBD}_t + \alpha * {lbdema}_{n - 1}$$
$\alpha$ is our \textbf{smoothing factor}, ${LBD}_t$ is the LBD of the current learned clause,
and \\${lbdema}_{n - 1}$ is the previous lbdema. In our testing of AMA VSIDS, I used 
$\alpha = 0.1$. 

\section{VSIDS - Results}
The results for VSIDS vary wildly because there are so many different types beyond what 
we've explored, and there are so many different parameters that could affect the end 
results. For VSIDS, although there might look like there's a clear winner, we need to 
keep in mind that we're using two relatively small test sets, so drawing an actual
conclusion as to which VSIDS version is ``objectively'' better isn't credible. 
Although AMA VSIDS isn't the fastest from our benchmarks, it will be the version we 
continue with for the rest of our paper. As for code, given how easy it is to swap out 
each version of VSIDS with each other, the exact code used for each version won't be 
provided, as the total changes were small enough that following along shouldn't be overly
difficult.

Another thing to note is that at the moment, there was not much exploration for the 
time differences. In order to figure out why the speeds between each version was 
different, it would be helpful to analyze how many conflicts there were, as well as the
order of deciding variables. As of now, I am just looking at the results and just 
accepting them as they are, not because I lack curiosity or scientific rigor, but because 
this is the second night in a row where I barely got any sleep for no reason, and I 
am too tired to analyze this under a microscope.

\begin{FVerbatim}
On all 1000 equations in uf20-91:

AMA (Y = 0.75 or 0.99):     274ms (~0.27 seconds).

EMA (X = 1, Y = 0.95):      266ms (~0.27 seconds).
EMA (X = 1, Y = 0.8):       287ms (~0.29 seconds).
EMA (X = 1, Y = 0.5):       280ms (~0.28 seconds).

mVSIDS (X = 1, Y = 0.5):    279ms (~0.28 seconds).
mVSIDS (X = 1, Y = 0.8):    277ms (~0.28 seconds).
mVSIDS (X = 10, Y = 0.5):   269ms (~0.27 seconds).
mVSIDS (X = 10, Y = 0.8):   284ms (~0.28 seconds).

cVSIDS (X = 1, Y = 0.5):    260ms (~0.26 seconds).
cVSIDS (X = 1, Y = 0.8):    266ms (~0.27 seconds).
cVSIDS (X = 10, Y = 0.5):   270ms (~0.27 seconds).
cVSIDS (X = 10, Y = 0.8):   279ms (~0.29 seconds).

Improved CDCL:              284ms (~0.28 seconds).

On all 1000 equations in UUF50.218.1000:

AMA (Y = 0.75 or 0.99):     1,629ms (~1.63 seconds).

EMA (X = 1, Y = 0.95):      1,504ms (~1.50 seconds).
EMA (X = 1, Y = 0.8):       1,497ms (~1.50 seconds).
EMA (X = 1, Y = 0.5):       1,621ms (~1.62 seconds).

mVSIDS (X = 1, Y = 0.5):    1,532ms (~1.53 seconds).
mVSIDS (X = 1, Y = 0.8):    1,526ms (~1.53 seconds).
mVSIDS (X = 10, Y = 0.5):   1,420ms (~1.42 seconds).
mVSIDS (X = 10, Y = 0.8):   1,467ms (~1.47 seconds).

cVSIDS (X = 1, Y = 0.5):    1,680ms (~1.68 seconds).
cVSIDS (X = 1, Y = 0.8):    1,705ms (~1.71 seconds).
cVSIDS (X = 10, Y = 0.5):   1,782ms (~1.78 seconds).
cVSIDS (X = 10, Y = 0.8):   1,841ms (~1.84 seconds).

Improved CDCL:              2,385ms (~2.39 seconds).
\end{FVerbatim}
\noindent